\section{Animal Classifier}
\label{sec:animal_classifier}

\subsection*{Learnings \& Troubleshooting}

\textbf{Problem: Falsche Bildauflösung (Resolution Shift)} \\
Wenn ein Convolutional Neural Network (wie \texttt{efficientnet\_v2\_s}) auf kleinen $224 \times 224$~Pixel großen Bildausschnitten trainiert wird, lernt es Filter für Objekte in einer bestimmten relativen Größe. Werden bei der Evaluierung oder bei der Visualisierung mittels Grad-CAM plötzlich $384 \times 384$~Bilder übergeben, kommt es zu einem sogenannten \emph{Resolution Shift}. Die gelernten Filter passen nicht mehr optimal zu den Strukturen, was die Accuracy künstlich verschlechtern kann.

\paragraph{Empfehlung:} 
Passe die Transformationen in der Datei \texttt{animalClassifier.py} so an, dass sie exakt der Trainingsauflösung von $224 \times 224$~Pixeln entsprechen. Alternativ kann in \texttt{train.py} direkt auf $384 \times 384$~Pixel trainiert werden, was allerdings mehr GPU-Speicher benötigt.

Both breed classifiers use EfficientNetV2-S \cite{tan2021efficientnetv2} as
the backbone, the small member of a model family designed jointly for
accuracy and training efficiency through a search over Fused-MBConv and
MBConv blocks. The choice was made in the preliminary phase and survived the
project unchanged: at ${\sim}21$M parameters the model is large enough for
fine-grained texture discrimination---breed identity frequently hinges on
coat texture, ear shape, and facial geometry rather than global
silhouette---yet small enough to train within a free-tier GPU budget and to
run comfortably inside the inference envelope: the profiling of
Sec.~\ref{sec:speed} attributes only $24.64$\,ms ($12.3\%$) of per-image
latency to the classification stage, an order of magnitude below the
detector, so the classifier is effectively free at inference time and the
two-model design carries no meaningful runtime penalty.

We instantiate the torchvision implementation with ImageNet-pretrained
weights---transfer learning being decisive at our dataset scale
\cite{kornblith2019transfer}, as the controlled pretrained-versus-scratch
comparison in Sec.~\ref{sec:training} quantifies---and replace the final
classification layer with a freshly initialized linear head sized to the
active class set: 10 outputs for the cat model, 10 for the dog model, and 20
in the earlier single-model edition that preceded the species split. All remaining layers are fine-tuned end-to-end rather than frozen: with tens
of thousands of images available (Sec.~\ref{sec:dataset}), the low-data
regime that motivates backbone freezing does not apply, and fine-grained
discrimination requires adapting precisely the deep texture features that
freezing would fix in place.

We additionally investigated whether the backbone initialization itself could
be improved for the harder species. Since ImageNet-1k devotes a large share
of its label space to dog breeds but only a handful of classes to cats, the
pretrained backbone arrives better prepared for the dog task than for the cat
task. To symmetrize this, we evaluated a community-published cat-breed
EfficientNetV2-S checkpoint from HuggingFace as an alternative cat-model
initialization. The experiment is a clean negative result: after unwrapping
the checkpoint's serialization format the backbone loaded correctly, yet final cat-model quality was indistinguishable from the
ImageNet baseline. Consistent with the data-centric findings of
Sec.~\ref{sec:training}, what ultimately improved the cat model was not
initialization but targeted dataset expansion. We retain the loader in the
repository for reproducibility.
